\section{From formulas to finite shift systems}
\label{sec:systems}

The word series of Section~\ref{sec:words} send every nonzero polynomial
$f=\sum_w f_w x_w$ to a nonzero series $h_f=\sum_w f_w h_w$.
To bound the order of this series in terms of formula size, we first represent
the word coefficients $f_w$ by matrices of bounded dimension.  We use the standard passage
from formulas to path programs~\cite[Section~2.1, Lemma~2, Step~1]{RS2005}
and the resulting coefficient representation~\cite[Lemma~3.1]{OAInc}.
We give the construction over an arbitrary field, keeping track of scalar
edges and the order of multiplication.

An \emph{acyclic scalar/variable-labelled path program} over a field $F$ is a
finite directed acyclic graph, with parallel edges allowed, whose edges are
labelled by scalars in $F$ or by variables $x_1,\ldots,x_n$.  It has designated
source and sink vertices; these vertices may coincide and need not be first
and last in a topological ordering.  The value of a directed path is the
product of its edge labels in traversal order.  The value of the program is
the sum of these products over all paths from the designated source to the
designated sink.  When the two vertices coincide, the empty path contributes
$1$.  Acyclicity makes this a finite sum in
$F\langle x_1,\ldots,x_n\rangle$.

\begin{lemma}[Coefficient representation of a path program]
\label{lem:acyclic-representation}
Let $F$ be a field, and let a path program with $m\geq 1$ vertices have value
$f=\sum_w f_w x_w$.  There are a row $u\in F^{1\times m}$, a column
$v\in F^{m\times 1}$, and strictly upper triangular matrices
$B_1,\ldots,B_n\in\Mat_m(F)$ such that
\begin{equation}
\label{eq:coefficient-representation}
 f_w=uB_wv,
 \qquad
 B_{i_1\cdots i_\ell}=B_{i_1}\cdots B_{i_\ell},
 \qquad B_\eps=I_m.
\end{equation}
In particular, $B_w=0$ whenever $|w|\geq m$.
\end{lemma}

\begin{proof}
Number the vertices in a topological order, so every edge goes from a smaller
index to a larger one.  Matrix rows index starting vertices and columns index
ending vertices.  Let $(E_0)_{ab}$ be the sum of the scalar labels on edges
from $a$ to $b$.  Let $(X_i)_{ab}$ be the sum of one copy of $1_F$ for each
edge from $a$ to $b$ labelled $x_i$.  All these matrices are strictly upper
triangular.  Consequently
\[
 S=I_m+E_0+\cdots+E_0^{m-1}
\]
has as its $(a,b)$ entry the sum of the values of all scalar-only paths from
$a$ to $b$, including the empty path when $a=b$.

Write $\mathrm{src}$ and $\mathrm{snk}$ for the designated vertices, and let
$e_a$ denote the $a$th coordinate column.  A path whose sequence of variable
labels is $i_1\cdots i_\ell$ decomposes uniquely into scalar-only paths
separated by these variable edges.  Since scalar labels are central, summing
over these decompositions gives
\begin{equation}
\label{eq:scalar-path-elimination}
 f_{i_1\cdots i_\ell}
 =e_{\mathrm{src}}^{\mathsf T}
   S X_{i_1}S\cdots X_{i_\ell}S e_{\mathrm{snk}}.
\end{equation}
For the empty word the right-hand side is
$e_{\mathrm{src}}^{\mathsf T}S e_{\mathrm{snk}}$.
Taking
\[
 u=e_{\mathrm{src}}^{\mathsf T}S,
 \qquad B_i=X_iS,
 \qquad v=e_{\mathrm{snk}}
\]
proves~\eqref{eq:coefficient-representation}.  The product of a strictly
upper triangular matrix with an upper triangular matrix is strictly upper
triangular, so each $B_i$ has this property.  Any product of $m$ such
$m\times m$ matrices is zero.
\end{proof}

\begin{lemma}[Formula representation]
\label{lem:formula-representation}
Let $F$ be a field, and let $f$ be computed by a noncommutative formula of
size at most $s$, counting all vertices including scalar leaves.  Then $f$
is the value of an acyclic scalar/variable-labelled path program with at
most $2s$ vertices.  Its coefficients therefore have a
representation~\eqref{eq:coefficient-representation} with $m\leq 2s$.
\end{lemma}

\begin{proof}
Associate two vertices, an entrance and an exit, to each formula vertex.
For a leaf, join its entrance to its exit by one edge labelled by the leaf's
variable or scalar.  At a sum gate, join its entrance to both child entrances
and both child exits to its exit, using edges labelled $1$.  At an ordered
product gate, use edges labelled $1$ from its entrance to the first child
entrance, from the first child exit to the second child entrance, and from
the second child exit to its exit.

The child subtrees have disjoint vertices because the computation is a
formula.  Recursively order the vertices of each subtree with its entrance
first, its first child subtree before its second, and its exit last.  Every
edge points forward in this ordering, so the graph is acyclic.  A path through
a sum gate chooses one of its two child programs; a path through a product
gate traverses the two child programs in their prescribed order.  Induction
on the formula therefore shows that the program from the root entrance to
the root exit has value $f$.  There are two program vertices per formula
vertex, including each scalar leaf.  The claimed coefficient representation
now follows from Lemma~\ref{lem:acyclic-representation}.
\end{proof}

We now return to a field $F$ of characteristic $p$ and to $K=F(t)$.
We combine the word series with the coefficient matrices
in~\eqref{eq:coefficient-representation} to form one invertible matrix.
Strict upper triangularity makes the following sum finite.

\begin{proposition}[Shift system for a coefficient representation]
\label{prop:formula-system}
Let $F$ be a field of characteristic $p$ and put $K=F(t)$.
Suppose $u\in F^{1\times m}$, $v\in F^{m\times 1}$, and strictly upper triangular
$B_1,\ldots,B_n\in\Mat_m(F)$ represent the coefficients of a polynomial $f$
as in~\eqref{eq:coefficient-representation}.  With the word series from
Section~\ref{sec:words}, define
\begin{equation}
\label{eq:fundamental-series-matrix}
 H(z)=\sum_w B_w h_w(z).
\end{equation}
Then $H\in\GL_m(K[[z]])$, $H(0)=I_m$, and
\begin{equation}
\label{eq:formula-shift-system}
 \begin{gathered}
 h_f=uHv,\qquad \sigma(H)=A(z)H,\\
 A(z)=I_m+\sum_{i=1}^n b_i(z)B_i=\frac{P(z)}{Q(z)},
 \end{gathered}
\end{equation}
where
\begin{equation}
\label{eq:system-denominator}
 \begin{gathered}
 Q(z)=\prod_{i=1}^n(1-a_i z),\qquad Q(0)=1,\\
 \deg_z Q\leq n,\qquad \deg_z P_{ab}\leq n
 \quad(1\leq a,b\leq m).
 \end{gathered}
\end{equation}
Here $P\in\Mat_m(K[z])$, and $a_i=(1+t)^i$ and
$b_i=z/(1-a_i z)$ are as in Section~\ref{sec:words}.
If $f\neq 0$, then $uHv\neq 0$.
\end{proposition}

\begin{proof}
Only words of length less than $m$ contribute to
\eqref{eq:fundamental-series-matrix}, so the sum is finite.
The empty word contributes $I_m$, whereas every nonempty word series has
zero constant term.  Hence $H(0)=I_m$ and $H$ is invertible over $K[[z]]$.
The coefficient representation gives
\[
 uHv=\sum_w (uB_wv)h_w=\sum_w f_w h_w=h_f.
\]
In particular, Proposition~\ref{thm:words} gives $uHv\neq 0$ when $f\neq 0$.

The shift fixes every $B_i$, and the defining recursion for the word series
gives $\sigma(h_{iw})-h_{iw}=b_i h_w$.  Splitting the nonempty words according
to their first letter, we obtain
\[
 \sigma(H)-H
 =\sum_{i=1}^n\sum_w B_iB_w b_i h_w
 =\left(\sum_{i=1}^n b_i B_i\right)H.
\]
This proves the shift equation with the displayed order of matrix factors.
Finally, multiplying $A$ by $Q$ gives
\[
 P(z)=Q(z)I_m+
       \sum_{i=1}^n z\prod_{j\neq i}(1-a_j z)B_i.
\]
Each summand has degree at most $n$, proving
\eqref{eq:system-denominator}.
\end{proof}

For a size-$s$ formula, Lemma~\ref{lem:formula-representation} makes the
system dimension at most $2s$.  The next section bounds the order of every
nonzero scalar component $uHv$ using this dimension and the degree bound in
\eqref{eq:system-denominator}.
